\chapter{Rendering, persistence, and diagnostics}\label{chap:output}
\section{A flat native 2D canvas}
\api{svg_view2<REAL>} maps the independent 2D world directly onto a rectangular
SVG image. Supply its world centre, minimum visible world width and height,
and optional image dimensions. The mapping uses one uniform scale; a viewport
with a different aspect ratio is letterboxed rather than stretching circles
into ellipses. Positive world $y$ points up while SVG image $y$ points down.

\begin{lstlisting}
namespace np = nurbspath;
using REAL = double;
const np::svg_view2<REAL> view({0.0, 0.0}, 6.0, 4.0, 900, 600);
np::svg_graphics_options2<REAL> style;
style.line_width = 1.5;
style.spline_segment_count = 160;
np::svg_canvas2<REAL> canvas(view, style);
canvas.add(np::point2<REAL>{-1.0, 1.0});
canvas.add(np::ray2<REAL>({-2.0, -1.0}, {1.0, 0.25}));
canvas.add(np::circle2<REAL>({0.5, 0.0}, 1.0));
const std::string svg = canvas.svg();
// canvas.write(output) streams the same self-contained document.
\end{lstlisting}
The canvas accepts native 2D points, rays, circles, and ordinary splines.
\api{add(spline)} draws its sampled polyline, control polygon, controls, and
evaluated endpoint dots. All flat-canvas strokes are solid, including
construction geometry. Unit red/green $XY$ references remain visible at the
2D origin. Vectors cannot be drawn alone because they do not supply a position.

\begin{figure}[htbp]\centering
\begin{tikzpicture}[x=1cm,y=1cm]
 \draw[surface] (-3,-1.7) rectangle (3,1.7);
 \draw[axis] (-2.7,-1.3)--(2.7,-1.3) node[below] {world $x$};
 \draw[axis] (-2.7,-1.3)--(-2.7,1.4) node[above] {world $y$};
 \draw[curve] (.2,.15) circle (.8);
 \node[point] at (.2,.15) {};
 \draw[manualred,thick] (-2.8,-.65)--(2.7,.65);
 \node[manualmuted] at (0,-2.05) {One scale for both world coordinates};
 \draw[manualorange,thick] (4.3,-1.7) rectangle (7.0,1.7);
 \draw[axis] (4.5,1.4)--(6.7,1.4) node[right] {SVG $x$};
 \draw[axis] (4.5,1.4)--(4.5,-1.15) node[below] {SVG $y$};
 \draw[curve] (5.65,0) circle (.7);
 \node[manualmuted,align=center] at (5.65,-2.05) {Image coordinates\\reverse $y$};
 \draw[-{Stealth},manualmuted,thick] (3.2,0)--(4.0,0);
\end{tikzpicture}
\caption{A flat drawing needs no plane embedding. Uniform scale preserves
shape; the image coordinate direction is a display convention.}
\label{fig:flat-canvas}
\end{figure}

\section{World-space cameras in 3D}
\api{svg_view3<REAL>::orthographic} takes a world-space eyepoint and look-at
point, orthographic world width/height, and image dimensions.
\api{svg_view3<REAL>::perspective} instead takes the vertical field angle in
radians and a positive near distance. A camera must have a nonzero viewing
direction. Image dimensions and orthographic extents must be positive; a
perspective field angle must be strictly between zero and $\pi$. Camera
getters expose the selected settings without changing the camera.

\begin{lstlisting}
const np::point3<REAL> eye{5.0, -7.0, 4.5};
const np::point3<REAL> target{0.0, 0.0, 0.5};
const auto ortho = np::svg_view3<REAL>::orthographic(
    eye, target, 6.0, 4.0, 900, 600);
const auto perspective = np::svg_view3<REAL>::perspective(
    eye, target, std::numbers::pi_v<REAL> / 4.0,
    900, 600, 0.05);
np::svg_document3<REAL> scene(perspective);
scene.add(np::sphere3<REAL>({0.0, 0.0, 0.0}, 1.0));
scene.add(np::plane3<REAL>({0.0, 0.0, -1.0}, {0.0, 0.0, 1.0}));
scene.add(np::ray3<REAL>({-2.0, 0.0, 0.0}, {2.0, 0.0, 0.0}));
const std::string scene_svg = scene.svg();
\end{lstlisting}

\begin{figure}[htbp]\centering
\begin{tikzpicture}[x=1cm,y=1cm]
 \node[manualink,font=\bfseries] at (1.9,2) {Orthographic};
 \draw[manualteal,thick] (0,-1.2)--(0,1.2);
 \draw[manualblue,thick] (3,-.8)--(3,.8);
 \foreach \y in {-.8,0,.8} \draw[-{Stealth},manualmuted] (3,\y)--(0,\y);
 \node[manualmuted] at (1.5,-1.6) {Equal size at every depth};
 \node[manualink,font=\bfseries] at (6.9,2) {Perspective};
 \node[point,label=below:Eye] (e) at (4.5,0) {};
 \draw[manualteal,thick] (5.5,-.7)--(5.5,.7);
 \draw[manualblue,thick] (8,-1.2)--(8,1.2);
 \draw[guide] (e)--(8,-1.2) (e)--(8,1.2);
 \draw[manualorange,very thick] (5.5,-.343)--(5.5,.343);
 \node[manualmuted,align=center] at (5.5,-1.25) {near/image\\plane};
 \node[manualmuted] at (7.1,-1.9) {Apparent size depends on depth};
\end{tikzpicture}
\caption{Camera projection is a display operation on the same 3D geometry.
The perspective diagram is a 2D cross-section of the viewing cone.}
\label{fig:camera-projection}
\end{figure}

Perspective world segments are clipped at the near plane before division by
depth. A segment crossing the plane retains its visible part; an entirely
hidden segment disappears. Orthographic output uses a uniform screen scale.
The sphere silhouette is sampled from geometry appropriate to the camera.

\begin{figure}[htbp]\centering
\begin{tikzpicture}
 \fill[manualred!6] (-3,-1.2) rectangle (0,1.5);
 \fill[manualteal!6] (0,-1.2) rectangle (4.2,1.5);
 \draw[manualorange,thick] (0,-1.2)--(0,1.5) node[above] {near plane};
 \draw[manualred,dashed,thick] (-2,-.6)--(0,.2);
 \draw[curve] (0,.2)--(3,1.4);
 \node[point] at (0,.2) {};
 \node[manualmuted] at (-1.7,1.1) {clipped};
 \node[manualmuted] at (2,-.7) {retained and projected};
 \draw[axis] (-2.7,-1.6)--(4,-1.6) node[right] {depth};
\end{tikzpicture}
\caption{Near-plane clipping prevents an invisible endpoint from producing
an unbounded perspective line. The clipped endpoint is an intermediate
rendering point, not a new library entity.}
\label{fig:near-clip}
\end{figure}

\section{Diagnostic glyphs and convenience rendering}
The default 3D scene is deliberately diagnostic rather than photorealistic.
The configurable \api{line_width} is in SVG viewBox units. The graphics options
also control spline/sphere segment counts, the plane square size, and other
documented construction lengths. \api{validate()} rejects invalid settings and
derived stroke sizes that cannot remain finite.

\begin{tabularx}{\linewidth}{@{}lY@{}}
\toprule
Entity & Default diagnostic conventions\\
\midrule
Point & Black location dot.\\
Ray & Pale-red dashed infinite supporting line, solid positive unit-direction
marker, and origin dot. Native direction magnitude does not change its glyph.\\
Sphere & Pale-blue silhouette, visible/hidden equator strokes, and centre dot.\\
Plane & Pale-green square (default side one quarter unit), unit normal,
half-unit positive $u$/$v$ markers, and parameter-origin dot.\\
Spline & Piecewise-linear curve, dashed control polygon, control dots, and
evaluated endpoint dots.\\
World axes & Persistent coloured $XYZ$ unit axes and labels at half line width.\\
\bottomrule
\end{tabularx}

A native 2D entity can be explicitly embedded for rendering with
\api{scene.add(plane, entity2)}. Its primary strokes are dash-dotted, with
dotted auxiliary/hidden/control strokes. The 2D circle bridge carries centre
and radius into a complete sphere. The separate flat canvas performs no such
embedding and always uses solid strokes.

\api{to_svg(view2, entity2, options2)} constructs a temporary flat canvas.
\api{to_svg(view3, entity3, options3)} constructs a 3D document.
\api{to_svg(view3, plane, entity2, options3)} performs explicit embedding. All
return self-contained XML strings; \api{svg()} and \api{write(output)} expose
the corresponding accumulated document. These functions do not render free
vectors or transform an arbitrary coordinate frame. A projected 2D circle
whose finite world radius overflows the screen radius has its outline omitted;
a finite centre still receives its dot.

\section{Coordinate streams, CSV/TSV, and binary records}
\api{point2}, \api{point3}, \api{vector2}, and \api{vector3} support
\api{operator<<} and \api{operator>>}. Their fields are whitespace-separated
coordinates in $x,y[,z]$ order. Reads commit only when a complete coordinate
record has been obtained. Vector/point semantics remain explicit even though
the scalar layout is similar.

\begin{lstlisting}
np::point3<REAL> position{1.0, 2.0, 3.0};
std::ostringstream csv;
position.csv_write(csv, np::text_format::csv, 4);
// Display: 1.0000e+00,2.0000e+00,3.0000e+00
std::istringstream csv_input(csv.str());
np::point3<REAL> restored;
restored.csv_read(csv_input, np::text_format::csv);
std::stringstream binary(std::ios::in | std::ios::out | std::ios::binary);
position.write(binary);
np::point3<REAL> binary_copy;
binary_copy.read(binary);
\end{lstlisting}
\api{csv_write(output, format, decimal_places)} uses scientific notation;
the default precision is six decimal places. Select \api{text_format::csv},
\api{tsv}, or \api{txt}. Delimited reads expect the chosen separators in the
coordinate record; this is not a general quoted-field CSV parser. These
helpers do not append a newline. Invalid format/negative output precision are
rejected. Stream errors use the normal stream state and exception mask.

Binary \api{write}/\api{read} store two or three native \api{REAL} values,
without a header or type discriminator. The layout is appropriate for matching
readers/writers, not portable interchange across byte orders, scalar types, or
floating-point representations. A failed complete read leaves the original
entity unchanged.

\section{Tagged records and independent presentation rows}
Tagged persistence supports four concrete 3D geometry types and three
presentation types. The scientific numbers in this displayed preview have
four decimal places; the canonical writers deliberately use full round-trip
\api{REAL} precision, independent of the destination stream formatting.

\begin{lstlisting}[style=plain]
42 point3 v1 1.0000e+00 2.0000e+00 3.0000e+00
42 vector3 v1 0.0000e+00 1.0000e+00 0.0000e+00
42 sphere3 v1 0.0000e+00 0.0000e+00 0.0000e+00 1.0000e+00
42 linecolor v1 2.0000e-01 4.0000e-01 8.0000e-01 7.5000e-01
42 linefont v1 1.5000e+00 4.0000e+00 2.0000e+00
42 text v1 1.2000e+01 "Path \"A\"\nnext pass"
\end{lstlisting}
\api{spline3} is the token for \api{nurbs_spline3}; its payload contains
degree, closure flag, tolerance, control and knot counts, interleaved control
$x,y,z,\text{weight}$ groups, then all knots. Array-backed splines convert to
an ordinary same-world spline before tagged output. There is no tagged 2D
geometry format or implicit 2D-to-3D persistence conversion.

\begin{tabularx}{\linewidth}{@{}lY@{}}
\toprule
Presentation value & Validation and meaning\\
\midrule
\api{line_color<REAL>} & \api{red}, \api{green}, \api{blue}, \api{alpha} are
finite in $[0,1]$. Alpha is opacity; default is opaque black.\\
\api{line_font<REAL>} & Finite positive width; \api{on}/\api{off} are both
zero for solid lines, or both finite positive for dashes. Default is
\api{{1,0,0}}.\\
\api{text_label<REAL>} & Finite positive size, default 12, and a string.
Width, dash lengths, and text size use consumer-selected display units.\\
\bottomrule
\end{tabularx}

\begin{lstlisting}
std::stringstream rows;
position.tag_write(42, rows);
np::line_color<REAL>{0.2, 0.4, 0.8, 0.75}.tag_write(42, rows);
np::line_font<REAL>{1.5, 4.0, 2.0}.tag_write(42, rows);
np::text_label<REAL>{12.0, "Path \"A\"\nnext pass"}.tag_write(42, rows);
while (auto record = np::tag_read_record<REAL>(rows)) {
    std::visit([tag = record->tag](const auto& owner) {
        // *owner has its concrete geometry or presentation type.
        (void)tag;
        (void)owner;
    }, record->value);
}
\end{lstlisting}

Tags are unsigned base-ten \api{std::size_t} application values, not unique
identifiers. Repeated tags, nonadjacent rows, and empty files are permitted.
The consumer chooses association, omitted-style defaults, and label placement.
Presentation rows do not alter SVG drawing state.

\api{T::tag_read(input)} returns
\api{std::optional<tagged_read_result<T>>}, with \api{tag} and a non-null
\api{std::shared_ptr<T>} named \api{entity}. A typed reader requires its
matching non-advisory type. \api{tag_read<REAL>} returns a
\api{tagged_entity3} whose \api{entity3_ptr} variant contains only point,
vector, sphere, and ordinary spline pointers. \api{tag_read_record<REAL>}
returns a \api{tagged_record} with the seven-type \api{value} variant.

Quoted labels escape backslash, double quote, newline, carriage return, and
tab. Each record remains one physical row. Other ASCII control bytes are
rejected; bytes at least 128 are preserved without UTF-8 validation. Writers
validate public fields before emitting invalid data.

\section{Advisory application and positioned messages}
\api{application_info} writes an untagged row with two nonempty, unquoted
tokens. ASCII whitespace/control bytes are not allowed in the tokens.
\api{diagnostic_message3<REAL>} selects \api{warning}, \api{information}, or
\api{error} using \api{diagnostic_level}, and supplies a finite 3D position,
positive finite size, and escaped text. The default is informational, at the
origin, size 12, with empty text.

\begin{lstlisting}
np::application_info{"path_exporter", "1.0"}.write(output);
np::diagnostic_message3<REAL>{
    np::diagnostic_level::warning, {1.0, 2.0, 3.0},
    12.0, "Small clearance"}.tag_write(42, output);
\end{lstlisting}
\begin{lstlisting}[style=plain]
application v1 path_exporter 1.0
42 warning v1 1.0000e+00 2.0000e+00 3.0000e+00 1.2000e+01 "Small clearance"
42 information v1 2.0000e+00 1.0000e+00 0.0000e+00 1.0000e+01 "Ready"
43 error v1 3.0000e+00 0.0000e+00 1.0000e+00 1.4000e+01 "No contact"
\end{lstlisting}
These are output-only annotations: readers may ignore them and the built-in
readers skip them. The application prefix alone identifies its advisory row;
a diagnostic must have a valid unsigned tag and a recognised diagnostic type.
Their versions and payloads are not decoded when skipped. Neither reader
variant gains an advisory alternative. An error annotation does not itself
terminate export or set the destination stream's failure state.

\begin{figure}[htbp]\centering
\begin{tikzpicture}[node distance=6mm]
 \node[draw=manualblue,fill=manuallight,rounded corners,minimum width=9.4cm,align=left] (a)
 {\api{application v1 tool 1.0}\quad\textcolor{manualmuted}{skip advisory}};
 \node[draw=manualblue,fill=manuallight,rounded corners,below=of a,minimum width=9.4cm,align=left] (b)
 {\api{42 warning v1 ...}\quad\textcolor{manualmuted}{skip advisory}};
 \node[draw=manualteal,fill=manualteal!8,rounded corners,below=of b,minimum width=9.4cm,align=left] (c)
 {\api{42 point3 v1 ...}\quad\textcolor{manualteal}{decode and return}};
 \node[draw=manualred,fill=manualred!5,rounded corners,below=of c,minimum width=9.4cm,align=left] (d)
 {\api{43 point3 v1 1 2}\quad\textcolor{manualred}{failbit: incomplete}};
 \draw[axis] (a)--(b); \draw[axis] (b)--(c); \draw[axis] (c)--(d);
\end{tikzpicture}
\caption{Reader calls skip advisory rows, then consume at most one required
geometry or presentation row. Malformed required data still fails strictly.}
\label{fig:tagged-reader}
\end{figure}

Malformed required rows set \api{failbit} after consuming their complete
physical row. Unknown types, unsupported required versions, missing/extra
fields, invalid tags, and invalid definitions fail. After clearing a
recoverable state, reading resumes from the next row; a reader never searches
past malformed required data. EOF and failure both return \api{nullopt}, so
inspect the stream state. CRLF and a final valid row without a newline are
accepted. Buffering and entity allocation can throw; applications needing
resource bounds must enforce their own line-size and complexity limits.

\section{Configuration and the hand-authored extension point}
\api{config.hpp} exposes semantic version strings/numbers and Git metadata,
including commit availability and dirty-state flags. Direct source includes
use the source fallback. A normal CMake build generates an updated header;
this standalone manual does not change that workflow. A repository without a
first commit can report an unborn state. Frequently used macros are:
\begin{itemize}
\item \api{NURBSPATH_VERSION_STRING}, \api{NURBSPATH_VERSION_NUMBER};
\item \api{NURBSPATH_GIT_COMMIT}, \api{NURBSPATH_GIT_COMMIT_AVAILABLE};
\item \api{NURBSPATH_GIT_DIRTY}, \api{NURBSPATH_GIT_DESCRIBE};
\item \api{NURBSPATH_CONFIG_GENERATED}.
\end{itemize}

The umbrella also includes \api{manual.hpp}, whose
\api{nurbspath::manual} namespace is an extension point for future hand-authored
functions. The current snapshot defines no algorithms there. It is distinct
from this document's top-level \api{manual/} directory.
